\begin{figure}[htbp]
			\centering
			\begin{tikzpicture}[
				scale=0.62,
				asta/.style={line width=1.5pt, line cap=round, line join=round, black},
				ridond/.style={line width=1.5pt, line cap=round, line join=round,
					gray!60},
				nodo/.style={circle, fill=white, draw=black,
					inner sep=0pt, minimum size=6pt, line width=1.2pt},
				quota/.style={thin, gray,  |<->|},
				richiam/.style={thin, gray, dashed},
				forza/.style={->, >=stealth, line width=1.5pt, black},
				]
				
				\def\SEP{11.0}
				
				%%── Pannello (a): cerniera in B ──────────────────────────────
				\begin{scope}[xshift=0cm]
					\coordinate (A) at (0,0); \coordinate (B) at (8,0);
					\coordinate (C) at (4,3); \coordinate (D) at (4,0);
					
					\draw[asta](A)--(C)node[midway,above left, font=\small]{$AC$};
					\draw[asta](C)--(B)node[midway,above right,font=\small]{$CB$};
					\draw[asta](A)--(D)node[midway,below,      font=\small]{$AD$};
					\draw[asta](D)--(B)node[midway,below,      font=\small]{$DB$};
					\draw[asta](C)--(D)node[midway,right,      font=\small]{$CD$};
					
					%% cerniera A
					\draw[line width=1.2pt](A)--++(-0.4,-0.65)--++(0.8,0)--cycle;
					\draw[line width=1.2pt](-0.55,-0.65)--(0.55,-0.65);
					\foreach \x in {-0.45,-0.25,-0.05,0.15,0.35}
					\draw[thin](\x,-0.65)--++(-.18,-.22);
					
					%% cerniera B
					\draw[line width=1.2pt](B)--++(-0.4,-0.65)--++(0.8,0)--cycle;
					\draw[line width=1.2pt]($(B)+(-0.55,-0.65)$)--($(B)+(0.55,-0.65)$);
					\foreach \x in {-0.45,-0.25,-0.05,0.15,0.35}
					\draw[thin]($(B)+(\x,-0.65)$)--++(-.18,-.22);
					
					\node[nodo]at(A){}; \node[nodo]at(B){};
					\node[nodo]at(C){}; \node[nodo]at(D){};
					
					\draw[forza]($(C)+(0,1.60)$)--($(C)+(0,0.35)$);
					\node[above,font=\normalsize]at($(C)+(0,1.60)$){$F$};
					
					\node[above left, font=\normalsize]at(A){$A$};
					\node[above right,font=\normalsize]at(B){$B$};
					\node[above right,font=\normalsize]at(C){$C$};
					\node[below=2pt,      font=\normalsize]at(D){$D$};
					
					\draw[richiam](A)--++(0,-1.3);
					\draw[richiam](D)--++(0,-1.3);
					\draw[richiam](B)--++(0,-1.3);
					\draw[quota](0,-1.1)--node[below,font=\small]{$\ell$}(4,-1.1);
					\draw[quota](4,-1.1)--node[below,font=\small]{$\ell$}(8,-1.1);
					\draw[richiam](C)--++(-4.3,0);
					\draw[richiam](A)--++(-0.8,0);
					\draw[quota](-0.6,0)--node[left,font=\small]{$h$}(-0.6,3);
					\node[font=\normalsize]at(4,-1.90){$(a)$};
				\end{scope}
				
				%%── Pannello (b): asta ridondante AB ─────────────────────────
				\begin{scope}[xshift=\SEP cm]
					\coordinate (A) at (0,0); \coordinate (B) at (8,0);
					\coordinate (C) at (4,3); \coordinate (D) at (4,0);
					
					\draw[asta](A)--(C)node[midway,above left, font=\small]{$AC$};
					\draw[asta](C)--(B)node[midway,above right,font=\small]{$CB$};
					\draw[asta](A)--(D)node[midway,above,      font=\small]{$AD$};
					\draw[asta](D)--(B)node[midway,above,      font=\small]{$DB$};
					\draw[asta](C)--(D)node[midway,right,      font=\small]{$CD$};
					
					%% asta ridondante AB (tratteggiata, leggermente sotto)
					\draw[ridond](A)--++(0,-0.28)--($(B)+(0,-0.28)$)--(B);
					\node[below,font=\small]at(4,-0.28){$AB$};
					
					
					%% cerniera A
					\draw[line width=1.2pt](A)--++(-0.4,-0.65)--++(0.8,0)--cycle;
					\draw[line width=1.2pt](-0.55,-0.65)--(0.55,-0.65);
					\foreach \x in {-0.45,-0.25,-0.05,0.15,0.35}
					\draw[thin](\x,-0.65)--++(-.18,-.22);
					
					
					%% carrello B
					\draw[line width=1.2pt](B)--++(-0.4,-0.50)--++(0.8,0)--cycle;
					\draw[line width=1.2pt]($(B)+(-0.50,-0.50)$)--($(B)+(0.50,-0.50)$);
					\draw[line width=1.2pt]($(B)+(-0.25,-0.65)$)circle(3.5pt);
					\draw[line width=1.2pt]($(B)+( 0.25,-0.65)$)circle(3.5pt);
					\draw[line width=1.2pt]($(B)+(-0.55,-0.82)$)--($(B)+(0.55,-0.82)$);
					\foreach \x in {-0.45,-0.25,-0.05,0.15,0.35}
					\draw[thin]($(B)+(\x,-0.82)$)--++(-.18,-.22);
					
					\node[nodo]at(A){}; \node[nodo]at(B){};
					\node[nodo]at(C){}; \node[nodo]at(D){};
					
					\draw[forza]($(C)+(0,1.60)$)--($(C)+(0,0.35)$);
					\node[above,font=\normalsize]at($(C)+(0,1.60)$){$F$};
					
					\node[above left, font=\normalsize]at(A){$A$};
					\node[above right,font=\normalsize]at(B){$B$};
					\node[above right,font=\normalsize]at(C){$C$};
					\node[above left,      font=\normalsize]at(D){$D$};
					
					\draw[richiam](A)--++(0,-1.3);
					\draw[richiam](D)--++(0,-1.3);
					\draw[richiam](B)--++(0,-1.3);
					\draw[quota](0,-1.1)--node[below,font=\small]{$\ell$}(4,-1.1);
					\draw[quota](4,-1.1)--node[below,font=\small]{$\ell$}(8,-1.1);
					\node[font=\normalsize]at(4,-1.90){$(b)$};
				\end{scope}
				
			\end{tikzpicture}
			\caption{Varianti iperstatiche della capriata triangolare.
				$(a)$~Iperstaticità esterna ($g_i^{\text{est}} = 1$):
				il carrello in~$B$ è sostituito da una cerniera;
				il sistema principale si ottiene ripristinando
				il carrello.
				$(b)$~Iperstaticità interna ($g_i^{\text{int}} = 1$):
				l'asta ridondante~$AB$ (in grigio) è aggiunta
				alla capriata isostatica; il sistema principale
				si ottiene sopprimendo~$AB$.}
			\label{fig:capriata_iperstatica}
		\end{figure}
